% Stage19 network-externalities theorem addendum (REVIEW-ONLY).
% Not connected to the canonical manuscript; no priority/journal claim.
\subsection{Network effects, residual design and the optimal depth of standards}

The network willingness-to-pay parameter $\eta$ is distinct from the
technical-substitutability weight $\rho$. Conditional demand is governed
by $B=\rho H+(1-\rho)X$, while fulfilled expectations obey
\[
 Mq=d\mathbf{1},\quad M=I+B-\eta H.
\]
Hence the net cross-firm Cournot interaction generated by an accepted
standards link of depth $g$ is $(\rho-\eta)g$, not $\rho g$.
Consumer surplus is $q^\top Bq/2$; an extra $\eta q^\top Hq$
term must NOT be added to welfare.

In a symmetric international-standard continuation define
\[
 \alpha=2-\eta,\quad\beta=1-\rho,\quad\delta=\rho-\eta,\quad
 t=\delta g+\beta y,\quad L=\alpha+2t,\quad P=\alpha-t
\]
and the marginal private gain from residual differentiation
\[
 A=\frac{6\beta\alpha d^2}{PL^3},\quad
 G=-A_t=\frac{6\beta\alpha d^2(5\alpha-8t)}{P^2L^4}.
\]

\begin{proposition}[Network--substitution balance in private product design]
Suppose the symmetric post-policy Nash design is interior and regular,
with $P>0$, $L>0$, $G>0$, and
$D=\kappa+\lambda-\beta G>0$. Then
\[
 \frac{\partial y^N}{\partial g}=
 \frac{(\rho-\eta)G}{D}.
\]
Consequently, deeper standards induce more private residual similarity
if $\rho>\eta$, do not change it if $\rho=\eta$, and induce private
residual differentiation if $\rho<\eta$.
The last case is explicitly conditional on a well-defined positive
quantity equilibrium, which is NOT globally assured across the original
strategy cube when $\eta>\rho$.
\end{proposition}
\begin{proof}
The single-firm unilateral Nash condition is
$\kappa(1-y)-\lambda(y-z)-A=0$.
Differentiate along a symmetric Nash continuation, using
$A_t=-G$ and $t_g=\delta+\beta y_g$.
Thus $-D y_g+\delta G=0$.
\end{proof}

\begin{proposition}[Network strength and strategic residual differentiation]
Hold $\rho,\kappa,\lambda,d$ and $g$ fixed and write
\[
 A_\eta=A\left[
 -\frac1\alpha+\frac{1-g}{\alpha-t}
 +\frac{3(1+2g)}{\alpha+2t}\right]>0.
\]
Whenever $t<\alpha$, $\alpha>0$, $P,L>0$ and $D>0$,
the fixed-installed-design comparative static satisfies
$y_\eta^N=-A_\eta/D<0$.
If the historical reference design is instead re-centered separately in
each compared economy at $z=z^{SW}(\eta)$, then
\[
 z_\eta=-\frac{A_\eta(z,0)}{\kappa-\beta G(\beta z)}<0,
 \qquad y_\eta^N=\frac{\lambda z_\eta-A_\eta(y,g)}{D}<0
\]
under the corresponding fragmentation curvature condition.
\end{proposition}
\begin{proof}
Since $t<\alpha$, $L<3\alpha$ and therefore
$3/L>1/\alpha$; also $(1-g)/P\ge0$.
The derivative signs then follow from implicit differentiation of the
private design Nash first-order conditions.
\end{proof}

\paragraph{Network and standards welfare decomposition.}
For a symmetric IS continuation with a genuinely fixed historical $z$,
let $N=\frac92+3(\rho g+\beta y)$ and
$W=N d^2/L^2-K(y)$. The fixed-design welfare effect of stronger network
willingness to pay is
\[
 W_\eta|_{g,y}=\frac{2Nd^2(1+2g)}{L^3}>0.
\]
When private Nash design is excessively similar so that
$W_y(y^N,g)<0$, strategic residual differentiation induced by stronger
network benefits adds $W_y y^N_\eta>0$. The network gain is reinforced
at a \emph{fixed} government depth. If, however, the legacy reference is
re-centered with $\eta$, the direct derivative acquires
$\lambda(y-z)z_\eta$ and no unconditional welfare sign is claimed.

\paragraph{Endogenous sovereign participation.}
The third-country joining surplus is
$\Phi(g,\eta)=W_C^{IS}(g,\eta)-W_C^{SU}(g,\eta)$, with each
regime evaluated at its respective private design Nash.
At a regular binding government optimum,
\[
 \frac{\partial g^*}{\partial\eta}
 =-\frac{\Phi_\eta}{\Phi_g},\qquad
 \frac{\partial y^*}{\partial\eta}
 =y^N_\eta+y^N_g\,g^*_\eta.
\]
The sign of this optimum-depth comparative static is NOT globally fixed
by the sign of $\rho-\eta$, because outsider participation incorporates
two different equilibrium designs and rent-shifting effects.
The model must be compared directly with Klimenko (2009) on
network-intensity-dependent over- and under-standardization.
